% s06-statement.tex —— §4 猜想的陈述（原文 4. Statement of the conjecture, p.28–29）
\section{猜想的陈述}% 原文 4. Statement of the conjecture
\srcpage{28}

与前面一样，$G$ 在 $E$ 上拟分裂，$\Gamma=\mathrm{Gal}(\overline{E}|E)$，并固定 $\ell\neq p$。
记 ${}^{L}G=\widehat{G}\rtimes\Gamma$ 为其 $\overline{\QQ}_\ell$-Langlands 对偶。
要注意，这个对偶应当通过猜想 3.21 内蕴地定义。
Langlands 参数按定义是连续态射
\[
\varphi:W_E\longrightarrow{}^{L}G,
\]
它到 $\Gamma$-因子的投影由嵌入 $W_E\subset\Gamma$ 给出。对这种参数记
\[
S_\varphi=\{g\in\widehat{G}\mid g\varphi g^{-1}=\varphi\}.
\]
这是某个代数群的 $\overline{\QQ}_\ell$-点集，该群的单位连通分支是约化群。总有
\[
Z(\widehat{G})^\Gamma\subset S_\varphi.
\]

\begin{definition}\label{def:4.1}
(1) 若 $S_\varphi/Z(\widehat{G})^\Gamma$ 有限，称 $\varphi$ 是离散的。

(2) 若 $\varphi$ 离散，且相应的 $1$-余圈 $I_E\to\widehat{G}$ 的像有限，称 $\varphi$ 是尖点的。
\end{definition}

通过假想的 Langlands 对应：
\begin{itemize}
\item 离散参数应当参数化离散序列的 $L$-包；
\item 尖点参数应当参数化超尖点的 $L$-包，即所有分量都超尖点的包。
\end{itemize}
在局部 Langlands 对应的所有已知情形里，这两条性质都成立。

\begin{remark}\label{rem:4.2}
在这个语境里并没有 tempered L-参数这个概念。事实上，这样的概念会依赖同构 $\overline{\QQ}_\ell\cong\CC$ 的选取。
若要求 $\varphi$ 下 Frobenius 的像的特征值都是 Weil 数，可以定义一个算术版本的概念；
但后一概念无法 $\ell$-进地插值，从纯粹局部的观点看也不自然。
\end{remark}

Langlands 参数构成群胚
\[
\mathcal{L}_G=\big[\mathrm{Hom}(W_E,{}^{L}G)/\widehat{G}\big].
\]
从这个观点看，离散参数对应这个“参数叠”成为 Deligne–Mumford 叠的那些点，差一个 $Z(\widehat{G})^\Gamma$-因子。

\begin{conjecture}[粗略版本]\label{conj:4.3}
存在群胚之间的函子
\[
\mathcal{L}_G^{\mathrm{disc}}\longrightarrow \mathrm{Perv}^{\mathrm{Weil}}(\Bun_G),\qquad
\varphi\longmapsto\cF_\varphi,
\]
满足以下性质：
\begin{enumerate}[label=(\arabic*)]
\item 该函子在 $\Bun_G$ 半稳定点处剩余 gerbe 上的茎，定义了 $G$ 的扩充纯内形式的局部 Langlands 对应；
\item 若 $\varphi$ 尖点，则 $\cF_\varphi$ 在 $\Bun_G$ 非半稳定轨迹上的限制为零；
\item $\cF_\varphi$ 是以 $\varphi$ 为特征值的 Hecke 本征向量。
\end{enumerate}
\end{conjecture}

下面给出更精确的表述。开始之前，必须先固定一个 Whittaker 数据。
事实上 $\cF_\varphi$ 的构造必须依赖这一选取：通过 $\cF_\varphi$ 构造出的 $G$ 的局部 $L$-包将有一个由 $S_\varphi$ 的平凡表示给出的特殊元素，
这个元素必须是该 $L$-包中唯一的 generic 元素。

\begin{conjecture}\label{conj:4.4}
给定离散 Langlands 参数 $\varphi$，存在 $\Bun_{G,\overline{\FF}_q}$ 上的“$\overline{\QQ}_\ell$-反常 Weil 层” $\cF_\varphi$，
带有 $S_\varphi$ 的作用，并满足以下性质：
\begin{enumerate}[label=(\arabic*)]
\item 对 $\alpha\in\pi_1(G)_\Gamma$，$Z(\widehat{G})^\Gamma$ 在 $\cF_\varphi$ 于分支 $\Bun_{G,\overline{\FF}_q}^\alpha$ 上限制的作用，
  由 $\alpha$ 通过等同 $\pi_1(G)_\Gamma=X^*(Z(\widehat{G})^\Gamma)$ 给出。
\item（尖点性条件）若 $\varphi$ 还是尖点的，则 $\cF_\varphi=j_!j^*\cF_\varphi$，其中 $j:\Bun_{G,\overline{\FF}_q}^{ss}\hookrightarrow\Bun_{G,\overline{\FF}_q}$。
\item（局部 Langlands 的实现）对基本的 $b\in G(L)$ 及 $x_b:[\Spa(\overline{\FF}_q)/J_b(E)]\hookrightarrow\Bun_{G,\overline{\FF}_q}$，
  作为 $S_\varphi\times J_b(E)$ 的表示、在 $J_b(E)$-分量上光滑，
  \[
  x_b^*\cF_\varphi=\bigoplus_{\substack{\rho\in S_\varphi\\ \rho|_{Z(\widehat{G})^\Gamma}=\kappa(b)}}\rho\otimes\pi_{\varphi,b,\rho},
  \]
  其中 $\{\pi_{\varphi,b,\rho}\}_\rho$ 是定义 $G$ 的扩充纯内形式 $J_b$ 的局部 Langlands 对应的 $L$-包。
  当 $b=1$ 时，$\pi_{\varphi,1,1}$ 是该 $L$-包中与 Whittaker 数据的选取相联系的唯一 generic 元素。
\item（Hecke 特征层性质）对 $\mu\in X_*(T)^+/\Gamma$，存在同构
  \[
  \overleftarrow{h}_!\big(\overleftarrow{h}^*\cF_\varphi\otimes IC_\mu\big)\cong\cF_\varphi\otimes r_\mu\circ\varphi.
  \]
  这里 $r_\mu\circ\varphi$ 是 $W_E$ 的 $\ell$-进表示，它在 $\Spa(E)^\diamond\times\Spa(\overline{\FF}_q)=\Spa(\breve E^{\mathrm{un}})^\diamond$ 上定义 Weil-étale 局部系统。
  该同构与 $S_\varphi$ 的作用相容，其中 $S_\varphi$ 在 $r_\mu\circ\varphi$ 上的作用由 $\varphi$ 诱导。
\item（特征层性质）对椭圆元 $\delta\in G(E)$，把它看作 $G(L)$ 的元素，Frobenius 在 $x_\delta^*\cF_\varphi$ 上的作用与 $\delta\in J_\delta(E)$ 的作用一致。
\item（局部/整体相容）设 $(H,X)$ 是 Hodge 型 Shimura 数据，$H_{\QQ_p}=G$，$\Pi$ 是 $H$ 的自守表示，使得 $\varphi_{\Pi_p}=\varphi$。
  则 Caraiani–Scholze 层 $R\pi_{\mathrm{HT}*}\overline{\QQ}_\ell[\dim\mathrm{Sh}]$（[2]）的 $\Pi^p$-等型分量，
  与 $\cF_\varphi$ 在 Hodge–Tate 周期 Grassmann 簇上的限制的某个倍数之间存在相容性。\srcpage{29}
\end{enumerate}
\end{conjecture}

我们稍后会把 (5) 与 (6) 说得更精确。这里指出：通过 Beauville–Laszlo 态射
\[
BL:\Gr\longrightarrow\Bun_G,
\]
“反常层”$\cF_\varphi$ 应当对应 $\Gr$ 上的一个“反常层”。但与 $IC_\mu$ 以及几何 Satake 同构中出现的其他反常层不同，
这个层复形在开 Schubert 胞腔上不会是局部常值的，其性质要复杂得多。

\begin{remark}\label{rem:4.5}
若 $G$ 的中心连通，则 $G$ 的所有内形式都是扩充纯内形式，即形如 $J_b$、$b$ 基本。
但一般情形下这不对，例如 $\mathrm{SL}_n$ 唯一的扩充纯内形式就是 $\mathrm{SL}_n$ 自身。
对 $\mathrm{SL}_n$，到达所有内形式的障碍位于 Galois 上同调群 $H^2(E,\mu_n)$，由局部类域论的基本类给出。
受 [6] 命题 3.4 的启发（该命题说这个基本类对应 $\oo(1)$ 的第一 Chern 类），可以把前述猜想扩充到 $\mathrm{SL}_n$ 的所有内形式：
改为看 $\oo(1)$ 的 $n$ 次方根 gerbe 上的 $\mathrm{SL}_n$-丛叠。作者希望类似的办法能推广到中心非连通的所有约化群 $G$。
\end{remark}

\begin{remark}\label{rem:4.6}
若在某个非基本的 $b$ 处取纤维 $x_b^*\cF_\varphi$，得到的是 $\widehat{J}_b$ 的表示，它必须经由 $\pi_0(\widehat{J}_b)=J_b(E)$ 分解（$\ell$ 与 $p$ 的不相容性）。
看来合乎逻辑的是：$J_b(E)$ 的这个光滑表示应当与某个联系于 Levi 子群 $M_b$（斜率态射的中心化子）的 Jacquet 函子相关（进一步评论见注 7.3）。
\end{remark}
